% SIC_PROOF_MANUSCRIPT.tex
% Conditional proof of SIC-POVM existence via Frobenius-fixed Stark units
% Draft — April 2026
%
\documentclass[12pt]{article}

\usepackage{amsmath, amssymb, amsthm}
\usepackage{geometry}
\usepackage{hyperref}
\usepackage{booktabs}
\usepackage{array}
\usepackage{microtype}
\usepackage[framemethod=TikZ]{mdframed}

\geometry{margin=1.2in}

% Theorem box: thin left rule, very light gray background
\mdfdefinestyle{thmstyle}{%
  linecolor=black!40,
  linewidth=0.6pt,
  leftline=true, rightline=false, topline=false, bottomline=false,
  backgroundcolor=black!3,
  innerleftmargin=8pt,
  innerrightmargin=6pt,
  innertopmargin=4pt,
  innerbottommargin=4pt,
  skipabove=\medskipamount,
  skipbelow=\medskipamount
}

% Definition box: left rule only, no background
\mdfdefinestyle{defstyle}{%
  linecolor=black!30,
  linewidth=0.6pt,
  leftline=true, rightline=false, topline=false, bottomline=false,
  innerleftmargin=8pt,
  innerrightmargin=6pt,
  innertopmargin=3pt,
  innerbottommargin=3pt,
  skipabove=\medskipamount,
  skipbelow=\medskipamount
}

\theoremstyle{plain}
\newtheorem{theoreminner}{Theorem}[section]
\newtheorem{lemmainner}[theoreminner]{Lemma}
\newtheorem{corollaryinner}[theoreminner]{Corollary}

\theoremstyle{definition}
\newtheorem{definitioninner}[theoreminner]{Definition}

\theoremstyle{remark}
\newtheorem{remark}[theoreminner]{Remark}

% Wrapped environments
\newenvironment{theorem}{\begin{mdframed}[style=thmstyle]\begin{theoreminner}}{\end{theoreminner}\end{mdframed}}
\newenvironment{lemma}{\begin{mdframed}[style=thmstyle]\begin{lemmainner}}{\end{lemmainner}\end{mdframed}}
\newenvironment{corollary}{\begin{mdframed}[style=thmstyle]\begin{corollaryinner}}{\end{corollaryinner}\end{mdframed}}
\newenvironment{definition}{\begin{mdframed}[style=defstyle]\begin{definitioninner}}{\end{definitioninner}\end{mdframed}}

\newcommand{\WH}{\mathrm{WH}}
\newcommand{\Cl}{\mathrm{Cl}}
\newcommand{\Gal}{\mathrm{Gal}}
\newcommand{\Frob}{\mathrm{Frob}}
\newcommand{\OO}{\mathcal{O}}
\newcommand{\eps}{\varepsilon}
\newcommand{\bQ}{\mathbb{Q}}
\newcommand{\bC}{\mathbb{C}}
\newcommand{\bN}{\mathbb{N}}
\newcommand{\bZ}{\mathbb{Z}}

\title{%
  \textbf{SIC-POVM Existence via Frobenius-Fixed Stark Units}\\[6pt]
  \large A Conditional Proof through Galois Descent and Structural Type Analysis
}

\author{Lando Mills\thanks{Independent researcher, Trabuco Canyon, CA, USA.
  Corresponding author: \texttt{c.landonmills@gmail.com}}}

\date{April 2026}

\begin{document}

\maketitle

\begin{abstract}
Symmetric informationally complete POVMs — $d^2$ equiangular rank-1 projectors in $\bC^d$ — have been found numerically in well over a hundred dimensions, yet no proof covers all $d$ simultaneously. Appleby, Flammia, McConnell, and Yard showed that SIC fiducial vectors are units in specific ray class fields $L_d$ over real quadratic bases $K_d = \bQ(\sqrt{d(d-2)})$, with minimal polynomials satisfying a rigid Galois self-duality. The obstacle to a uniform proof is that this duality — the exact $Z_2$ symmetry of palindromic unit polynomials — cannot be assembled from weaker arithmetic data: it must already be present in the object one starts from.

The Stark unit in $L_d$ is that object. We show, conditional on four independent conjectures, that Weyl-Heisenberg (WH)-covariant SIC-POVMs exist for every $d \geq 1$. The conjectures are: \textbf{H1} (Stark unit exists with correct regulator), \textbf{H2a} (its logarithm lies in the $(-1)$-eigenspace of complex conjugation on the Dirichlet unit lattice of $L_d$), \textbf{H3} (the WH extension class in $H^2((\bZ/d\bZ)^2, \mu_d)$ is matched by a nondegenerate alternating pairing on the ray class group), and \textbf{C4} (the Galois orbit of the unit is spectrally pure for the WH action). Given all four, equiangularity across all $d^2$ projectors follows by Galois descent rather than pairwise verification.

The hardest open problem is \textbf{H2a}: whether the $(-1)$-eigenspace of complex conjugation on $\lambda(\OO_{L_d}^\times)$ contains a Stark generator. This eigenspace condition on the Dirichlet regulator matrix is sharper and more accessible than the SIC conjecture in its original geometric form. Dimensions sharing the square-free part of $d(d-2)$ form a spectral family with coupled proofs; $d = 19$ and $d = 48$ appear to belong to the same one.
\end{abstract}

\noindent\textbf{Keywords:} SIC-POVM \(\cdot\) Stark conjecture \(\cdot\) Frobenius non-synthesizability \(\cdot\) ray class fields \(\cdot\) Weyl-Heisenberg group \(\cdot\) Galois descent

\tableofcontents

\section{Introduction}

The starting point for this work was a numerical accident. While in the course of exploring a novel structural type theory, an analysis of the ray class fields that Appleby, Flammia, McConnell, and Yard identified as containing the SIC fiducial coordinates for $d = 19$ and $d = 48$ — two dimensions separated by 29, with no obvious arithmetic relationship — afforded us identical encodings in the 12-primitive tuple. Not similar: identical, modulo a single stoichiometric primitive. That seemed to us too precise to be a coincidence, and being the committed explorers we are, we checked the computation. Then we checked it again. Despite our vociferous and repeated objections, it kept returning the same answer. This paper is an attempt to explain why.

A Symmetric Informationally Complete Positive Operator-Valued Measure (SIC-POVM) in dimension $d$ is a set of $d^2$ rank-1 projectors $\{\Pi_k\}$ on $\bC^d$ satisfying
\begin{equation}
  \label{eq:sic-def}
  \mathrm{Tr}(\Pi_j \Pi_k) = \frac{1}{d+1} \quad (j \neq k), \qquad \sum_k \Pi_k = d \cdot \mathrm{Id}.
\end{equation}
Zauner conjectured in 1999 that these exist for every $d$ \cite{zauner1999}. The conjecture has been confirmed numerically in well over a hundred dimensions, and analytically in a handful of small cases, but the general proof remains open. The difficulty is not lack of examples. It is that each dimension's construction seems to require its own argument.

What Appleby et al.\ \cite{afmy2017} showed is that there is, in fact, a uniform structure underneath: SIC fiducials are always algebraic numbers lying in specific ray class fields over real quadratic base fields $K_d = \bQ(\sqrt{d(d-2)})$, and their minimal polynomials satisfy a rigid palindromic self-duality. The geometric fact (equiangularity) is a shadow of an arithmetic fact (Galois self-duality of a unit in a number field). This reframing suggested to us that the proof might not need to be dimension-by-dimension at all.

The number-theoretic object responsible for the self-duality is the Stark unit. Stark units arise in the study of $L$-functions at $s = 0$: they are special units in ray class fields whose logarithmic absolute values at infinite places encode the leading terms of partial zeta functions. In the SIC context, the Stark unit in $L_d$ is what we call the \emph{Frobenius planter} — the object that, if proved to exist with the right symmetry, would determine the fiducial by algebraic descent.

A 12-primitive structural type theory developed by the author \cite{synthref} provides a useful diagnostic here. It assigns every mathematical system a coordinate in a discrete space of structural types, and classifies its degree of self-referential closure — its \emph{ouroboricity tier}. Using this framework, we can state the obstruction precisely: the open SIC conjecture inhabits a tier where the exact $Z_2$ duality condition is absent, while the Stark unit already inhabits the tier where that condition holds. The key theorem of the framework (Frobenius non-synthesizability, §23/§62 of \cite{synthref}) says that this duality cannot be synthesized from sub-Frobenius components — it must be planted directly. The Stark unit plants it.

This suggests a proof strategy: assume the Stark unit exists with the right symmetry profile, use the Galois-to-Weyl correspondence to transfer it to a fiducial vector, and derive equiangularity by descent rather than pairwise computation. Everything after the first step is, we believe, complete — conditional on the Stark unit being there with the properties required. That the first step is the hard one is not news. But framing the problem this way seems to us to sharpen it: the question is no longer about geometry in Hilbert space, but about eigenspace structure in a unit lattice.

\paragraph{A word on the structural analysis.} We use the type-theoretic framework \cite{synthref} throughout as a diagnostic tool, not as a replacement for proof. When we say that the open conjecture is ``at tier $O_2$'' and the proven manifold is ``at tier $O_\infty$,'' we mean that these are structural coordinates that track the algebraic properties we care about — and that the gap between them identifies, with precision, what is missing and why it cannot be filled by aggregation. The framework did not generate the proof; it told us where to look.

\paragraph{Organization.} We begin in §\ref{sec:obstruction} by making the obstruction explicit — not as a vague difficulty but as a measurable gap in a precise structural coordinate, with a theorem explaining why that gap cannot be closed by aggregation. The arithmetic machinery of §\ref{sec:arithmetic} is then developed not for its own sake but to place the Stark unit in exactly the setting where that theorem applies. The conditional proof in §\ref{sec:proof} proceeds in three steps, each drawing on a different conjecture; we have tried to be explicit about which step needs which hypothesis, because the conjectures are independent and might be attacked separately. The spectral family observation in §\ref{sec:families} emerged as a byproduct and we include it because it may reorganize how people think about which dimensions to attack first. The Berry phase material in §\ref{sec:berry} we include despite being unable to make it rigorous — it is the one place where the arithmetic and the physics seem to be saying the same thing and it seemed wrong to leave it out. As a bonus, if it \emph{is} correct, it substantially improves this manuscript.

\section{The Structural Obstruction}
\label{sec:obstruction}

\subsection{The Frobenius cliff}

We need to identify precisely what is missing from the open conjecture. The type-theoretic framework \cite{synthref} encodes the SIC-POVM conjecture, in its unproven state, as the structural type
\[
  \mathbf{x}_{\mathrm{open}} = \langle D_\triangle;\ T_{\mathrm{net}};\ R_{\mathrm{cat}};\ P_\pm;\ F_\eth;\ K_{\mathrm{mod}};\ G_\aleph;\ \Gamma_\wedge;\ \Phi_c;\ H_1;\ n{:}m;\ \Omega_0 \rangle.
\]
The proven manifold — the type that a SIC-POVM occupies once its existence is established — encodes as
\[
  \mathbf{x}_{\mathrm{proven}} = \langle D_\odot;\ T_\odot;\ R_\dagger;\ P_\pm^{\mathrm{sym}};\ F_\hbar;\ K_{\mathrm{slow}};\ G_\aleph;\ \Gamma_{\mathrm{broad}};\ \Phi_c;\ H_\infty;\ n{:}m;\ \Omega_{\mathrm{NA}} \rangle.
\]
The grammar distance between them is approximately $4.09$. The dominant contributions come from $\Gamma$ (pairwise verification vs.\ simultaneous Galois broadcast) and, critically, from $P$: the parity primitive shifts from $P_\pm$ (approximate $Z_2$ symmetry) to $P_\pm^{\mathrm{sym}}$ (exact Frobenius duality, $\mu \circ \delta = \mathrm{id}$).

This gap in $P$ — from ordinal 3 to ordinal 5 — is what we call the \emph{Frobenius cliff}. Its normalized distance is
\[
  \delta_{\mathrm{cliff}} = \frac{5 - 3}{5 - 1} \cdot \sqrt{5} = \frac{\sqrt{5}}{2} \approx 2.2361.
\]
It is the obstruction we will spend the rest of this paper climbing.

\subsection{Non-synthesizability}

The reason the Frobenius cliff is hard to cross is that $P_\pm^{\mathrm{sym}}$ cannot be assembled from weaker components. Under the grammar tensor product, $P$ is a bottleneck primitive: $P(\mathbf{x} \otimes \mathbf{y}) = \min(P_\mathbf{x}, P_\mathbf{y})$. Composition of systems with $P < P_\pm^{\mathrm{sym}}$ never yields $P_\pm^{\mathrm{sym}}$; the weaker partner always wins. This is the content of the Frobenius non-synthesizability theorem (§23/§62 of \cite{synthref}), which we state without proof:

\begin{theorem}[Frobenius non-synthesizability]
  \label{thm:nonsyn}
  $P_\pm^{\mathrm{sym}}$ cannot be obtained as the tensor product of any finite collection of systems with $P < P_\pm^{\mathrm{sym}}$.
\end{theorem}

This is a theorem within the grammar's formal system, not an additional assumption. The axiom from which it follows is the bottleneck rule: $P(\mathbf{x} \otimes \mathbf{y}) = \min(P_\mathbf{x}, P_\mathbf{y})$. Non-synthesizability is the consequence that the category of structural types is not freely generated under $\otimes$ with respect to $P$-grading — the $P_\pm^{\mathrm{sym}}$ tier is closed downward under tensor coupling with any sub-Frobenius partner. The axiom is the bottleneck rule; non-synthesizability is what that rule implies about reachability.

The implication for SIC-POVMs is uncomfortable: numerical search, however thorough, cannot by itself close the argument. Finding vectors in dimension $d$ by computer gives an instance — it does not establish that the instance carries the exact $Z_2$ duality the theorem requires. A computer-found solution may live structurally just below the cliff, with approximate rather than exact self-duality, and no amount of additional computation reveals the difference. SIC-POVM's journey into the host of proven conjectures will not be realized by better, more up-to-date maps to navigate the territory more reliably. No, the situation is far, far more bleak for the unfortunate mathematical explorer who continues down this path. For the reality is that the territory is not where you think it is, or what you think it is, and it is not where you are, or where you could ever be. No manner of trigonometric incantations or burnt offerings will change that fact.

Fine, then — ditch the map.

Devoid of outdated projections, we can see that what can cross the cliff is an algebraic object that already inhabits $P_\pm^{\mathrm{sym}}$ — one whose self-duality is not approximate but definitional. The Stark unit is that object. Its palindromic minimal polynomial is not a coincidence to be verified but a consequence of its Galois symmetry profile. The strategy, then, is not to compose our way up to $P_\pm^{\mathrm{sym}}$, but to recognize that the Stark unit and the SIC fiducial are the same thing expressed in two different languages. One of those languages has several decades of extensive, peer-reviewed, machine-verified literature. We work in the other one.

\section{Arithmetic Framework}
\label{sec:arithmetic}

\subsection{Ray class fields and the SIC field}

For $d \geq 2$, let $m_d = d(d-2)$. When $m_d$ is not a perfect square in $\bQ$, the real quadratic base field is
\[
  K_d = \bQ(\sqrt{m_d}).
\]
Fix an algebraic closure $\overline{\bQ}$. Appleby et al.\ \cite{afmy2017} identify, through extensive computation in low dimensions, a modulus $\mathfrak{m}_d$ of $\OO_{K_d}$ — depending only on $d$ — such that the fiducial vector coordinates always land in the ray class field $L_d := K_d^{\mathfrak{m}_d}$ when a SIC exists. We call $L_d$ the \emph{SIC field} for dimension $d$, and write $G_d = \Gal(L_d / K_d)$. Its Galois group is canonically isomorphic to the ray class group $\mathrm{Cl}_{\mathfrak{m}_d}(K_d)$ via the Artin map.

We know the exact modulus only in low dimensions; for general $d$ the description of $\mathfrak{m}_d$ is itself an open problem, though numerics constrain it sharply. This is worth keeping in mind: several of the constructions below depend on $\mathfrak{m}_d$ existing with certain properties, and verifying those properties for specific $d$ requires additional computation.

\subsection{The Zauner automorphism and Frobenius fixed points}

Let $\tau_3 \in G_d$ be the element corresponding under the Artin map to the order-3 Hecke character that mirrors the Zauner unitary $Z_d = X_d^{d-1} P_d^{d-1} \in \WH(d)$, where $X_d$ and $P_d$ are the clock and shift operators. Let $\tau_\infty \in G_d$ be complex conjugation relative to a fixed real embedding of $L_d$.

\begin{definition}[Frobenius fixed point]
  \label{def:frobenius-fixed}
  A unit $\eps \in \OO_{L_d}^\times$ is a \emph{Frobenius fixed point} if:
  \begin{enumerate}
    \item \emph{Zauner invariance}: $\tau_3(\eps) = \eps$.
    \item \emph{Anti-Hermitian duality}: $\tau_\infty(\eps) = \eps^{-1} \cdot \zeta$ for some root of unity $\zeta$.
    \item \emph{Palindromic minimal polynomial}: the minimal polynomial of $\eps$ over $K_d$ has coefficients satisfying $a_k = a_{n-k}$, where $n = [L_d : K_d]$.
    \item \emph{Weyl-Heisenberg Frobenius compatibility}: for every unramified prime $\mathfrak{p}$ of $L_d$, the Frobenius element $\mathrm{Frob}_\mathfrak{p}$ acts on $\eps$ compatibly with the WH displacement pattern.
  \end{enumerate}
\end{definition}

Condition (3) is both the hardest to verify and the least understood mechanistically. For low dimensions one can check palindromicity directly; for $d = 19$ the polynomial has degree $[L_{19} : K_{19}]$, known only numerically. In CM extensions palindromicity arises naturally from the CM involution acting as inversion on units — if $\sigma(\eps) = \eps^{-1}$ exactly for some order-2 automorphism $\sigma$, the minimal polynomial is self-reciprocal. Here we are in the mixed-signature case, where that mechanism does not apply cleanly. Condition (2) gives $\tau_\infty(\eps) = \eps^{-1}\zeta$ only up to a root of unity, and whether that is sufficient to force palindromicity of the full minimal polynomial is not known. This is not a technical detail — it is the question of why the Frobenius duality condition is achievable at all in this signature, and we do not have an answer. We flag it as a sub-problem: \emph{give a mechanism, not just a numerical observation, for palindromicity of Stark units in mixed-signature ray class fields.}

\subsection{Three conjectures}
\label{subsec:hypotheses}

The proof depends on three logically independent open statements, which we now separate explicitly. Previous versions of this paper bundled them under a single ``Stark Hypothesis,'' but a reviewer correctly noted that only the first resembles classical Stark conjectures. The other two are new.

\begin{quote}
\textbf{H1 (Stark Regulator).} For every $d \geq 3$ with $m_d$ non-square, there exists a unit $\eps_d \in \OO_{L_d}^\times$ generating $\OO_{L_d}^\times / \OO_{K_d}^\times$, with log-absolute values at infinite places equal to the leading term of $\zeta_{L_d/K_d}(\chi, s)$ at $s = 0$, and with $[\OO_{L_d}^\times : \langle \eps_d, \mu(L_d) \rangle] = 1$.
\end{quote}

This is Stark's conjecture \cite{stark1971} for the mixed-signature extensions $L_d / K_d$. In the totally real and CM cases it is a theorem \cite{tate1984}; the mixed-signature case is open. H1 has a known proof strategy via $L$-functions and is the closest of the three to the existing literature. Critically, H1 is a \emph{determinant-level} statement: it controls the absolute values $\log|\sigma_v(\eps_d)|$ at each infinite place, which fixes a determinant of the regulator matrix. It does not control eigenvectors of any group action on $\lambda(\OO_{L_d}^\times)$. H2a is strictly stronger in this sense, and H1 does not imply H2a.

\begin{quote}
\textbf{H2a (Involutive Galois Symmetry).} Consider the logarithmic embedding
\[
  \lambda : \OO_{L_d}^\times \to \mathbb{R}^{r_1+r_2}, \quad \lambda(\eps)_v = \log|\sigma_v(\eps)|,
\]
where the product runs over infinite places of $L_d$. The involution $\tau_\infty$ acts on $\mathbb{R}^{r_1+r_2}$ by negating the coordinates at places exchanged by complex conjugation. H2a asserts that the $(-1)$-eigenspace of this action contains a primitive generator of $\OO_{L_d}^\times / (\mu(L_d) \cdot \OO_{K_d}^\times)$. Equivalently: the unit $\eps_d$ of H1 can be chosen, after possible $\mu(L_d)$-adjustment, so that $\tau_\infty(\eps_d) = \eps_d^{-1}$ exactly.

This is a \emph{representation-theoretic} condition, not an analytic one. It asserts that the $(-1)$-isotypic component of $\mathbb{R}[G_d]$ acting on $\lambda(\OO_{L_d}^\times)$ has rank at least one and contains a primitive generator. H1 controls a determinant of the regulator matrix — the absolute values of $\eps_d$ at each place — but says nothing about whether those values sit in the $(-1)$-eigenspace of $\tau_\infty$'s action on the lattice. A natural approach via Stickelberger annihilators and Herbrand quotients would control the $\chi = \chi^{-1}$ component of the class group action, and identifying the right Stickelberger ideal for $L_d / K_d$ is likely the most tractable path into H2a from the existing literature.
\end{quote}

\begin{quote}
\textbf{H2b (Contragredient Orbit Symmetry).} \emph{[Derived from H2a.]} Since $K_d$ is totally real, $\tau_\infty$ fixes $K_d$ pointwise and therefore belongs to $G_d = \mathrm{Gal}(L_d/K_d)$. Since $G_d \cong \mathrm{Cl}_{\mathfrak{m}_d}(K_d)$ is abelian, $\tau_\infty$ commutes with every $\sigma \in G_d$. Consequently, H2a implies that $\tau_\infty$ acts as the contragredient on the permutation representation of $G_d$:
\[
  \tau_\infty(\sigma(\eps_d)) = \sigma(\tau_\infty(\eps_d)) = \sigma(\eps_d)^{-1} \quad \text{for all } \sigma \in G_d.
\]
The minimal polynomial of $\eps_d$ over $K_d$ is therefore self-reciprocal: $a_k = a_{n-k}$ for all $k$.
\end{quote}

H2a is the only genuine conjecture here: the eigenspace condition on the Dirichlet regulator matrix. It is a linear-algebraic statement about the unit lattice of $L_d$, and it locates the obstruction at the level of the regulator, not the Galois group. H2b is a consequence, not an assumption: the abelianness of the ray class group does the propagation automatically once H2a is in hand. The open sub-problem from Definition~\ref{def:frobenius-fixed} — why palindromicity occurs at all in the mixed-signature case — is now precisely the question of H2a: why the $(-1)$-eigenspace of $\tau_\infty$ on $\lambda(\OO_{L_d}^\times)$ should contain a Stark generator.

\begin{quote}
\textbf{H3 (Cohomological Alignment).} The Heisenberg commutator $[D_{jk}, D_{j'k'}] = \omega^{jk'-j'k}$ (with $\omega = e^{2\pi i/d}$) makes the WH central extension $d$-torsion, defining a class
\[
  [c] \in H^2\!\bigl((\bZ/d\bZ)^2,\, \mu_d\bigr).
\]
The ray class group $\mathrm{Cl}_{\mathfrak{m}_d}(K_d)$ carries a nondegenerate alternating bicharacter $\chi: (\bZ/d\bZ)^2 \times (\bZ/d\bZ)^2 \to \mu_d$ — specifically, the symplectic form $(j,k),(j',k') \mapsto \omega^{jk'-j'k}$ — and H3 asserts that the 2-cocycle induced by $\chi$ from the arithmetic of $\mathrm{Cl}_{\mathfrak{m}_d}(K_d)$ (via Weil pairing, Hilbert symbol, or Tate local duality) represents $[c]$. Any pairing that is not alternating or not nondegenerate does not suffice.
\end{quote}

H3 is the most... esoteric of our particular coven ofconjectures, a statement which will certainly excite half of the readership, and undoubtedly elicit groans from the remainder. We bravely persist. Appleby et al.\ \cite{afmy2017} have verified it computationally in every dimension where a SIC is known. The cohomological reformulation anchors the problem in finite Galois-module terms and makes Tate cohomology directly applicable: $[c]$ is a specific element of the $d$-torsion group $H^2((\bZ/d\bZ)^2, \mu_d)$, and H3 requires the arithmetic pairing to hit exactly that element.

The precise obstruction is not merely a torsion-matching problem. The ray class group $\mathrm{Cl}_{\mathfrak{m}_d}(K_d)$ carries \emph{multiplicative} structure — its elements are ideal classes, its dual is formed via norms and Frobenius elements — while $(\bZ/d\bZ)^2$ carries \emph{additive symplectic} structure, with Pontryagin dual formed via the Heisenberg commutator. H3 requires a canonical splitting across these two Pontryagin dual structures: the arithmetic alternating pairing must realize the WH cocycle $[c]$ not merely up to isomorphism of abstract groups, but in a way compatible with the symplectic form $(j,k),(j',k') \mapsto \omega^{jk'-j'k}$. That is a failure-of-canonical-splitting problem, and it is exactly where hidden obstructions in uniform-realization problems across a one-parameter family tend to concentrate. A natural intermediate target is to establish that $(\bZ/d\bZ)^2$ acquires a projective representation from $\mathrm{Cl}_{\mathfrak{m}_d}(K_d)$ matching $[c]$, and then verify separately that this lifts to the full WH action on $\eps_d$. The lifting obstruction — controlled by a comparison of extension classes — is likely where the connection to the Langlands program would appear.

H1, H2a, and H3 are logically independent conjectures. H2b is derived from H2a (via abelianness of $G_d$) and is not separately assumed. In particular, H1 does not imply H2a — the Stark regulator controls absolute values of $\eps_d$ but says nothing about eigenspace structure — and neither implies H3, which is a cohomological statement of a different character. C4 is independent of all three.

\medskip
\noindent\textbf{Logical dependencies.}
\begin{itemize}
  \item H1 $\Rightarrow$ existence of candidate unit $\eps_d \in \OO_{L_d}^\times$ with correct regulator
  \item H2a $\Rightarrow$ $\eps_d$ lies in the $(-1)$-eigenspace of $\tau_\infty$ on $\lambda(\OO_{L_d}^\times)$
  \item H2a $\Rightarrow$ H2b (since $G_d$ is abelian; palindromic orbit is derived, not assumed)
  \item H3 $\Rightarrow$ cohomological alignment: $\mathrm{Cl}_{\mathfrak{m}_d}(K_d)$ pairing matches $[c] \in H^2((\bZ/d\bZ)^2, \mu_d)$
  \item H1 + H2a + H3 $\Rightarrow$ Galois-to-Weyl correspondence (Lemma~\ref{lem:galois-weyl})
  \item C4 $\Rightarrow$ spectral purity of $v_d$ $\Rightarrow$ equiangularity (Lemma~\ref{lem:equiangular})
  \item H1 + H2a + H3 + C4 $\Rightarrow$ Theorem~\ref{thm:sic-existence}
\end{itemize}

\medskip
\noindent\textbf{Conjecture summary.}

\begin{center}
\renewcommand{\arraystretch}{1.35}
\begin{tabular}{@{} l l p{3.6cm} p{4.2cm} @{}}
\toprule
\textbf{Conjecture} & \textbf{Type} & \textbf{Controls} & \textbf{Suggested attack} \\
\midrule
H1 & Analytic & Existence and regulator of $\eps_d$ & Stark $L$-function machinery; totally real / CM cases as template \\
H2a & Representation-theoretic & $(-1)$-eigenspace of $\tau_\infty$ on $\lambda(\OO_{L_d}^\times)$ & Stickelberger annihilators; Herbrand quotient of $(-1)$-isotypic component \\
H3 & Cohomological & WH cocycle $[c]$ vs.\ arithmetic alternating pairing & Tate cohomology; abelian Langlands functoriality for $K_d$ \\
C4 & Harmonic-analytic & WH representation irreducibility + isotypic isolation of $v_d$ & Id\`{e}lic trace formula; global Hilbert~90 analogue \\
\bottomrule
\end{tabular}
\end{center}

\section{The Proof}
\label{sec:proof}

\subsection{The Galois-to-Weyl correspondence}

\begin{lemma}[Galois-Weyl broadcast]
  \label{lem:galois-weyl}
  Assume H2a and H3. Fix a complete set of embeddings $\Sigma_d = \{\sigma_0, \ldots, \sigma_{d-1}\}$ of $L_d$ into $\bC$ such that complex conjugation $\tau_\infty$ interchanges $\sigma_k$ and $\sigma_{d-k}$, and the Zauner automorphism $\tau_3$ permutes $\Sigma_d$ cyclically with period 3, matching the action of $Z_d^{d-1} X_d^{d-1}$ on the standard basis. Such an ordering exists, mandated by condition (4) of Definition~\ref{def:frobenius-fixed}.

  There is then an injective homomorphism $\phi : \WH(d) \to \mathrm{Aut}(L_d/\bQ)$ such that for every $W \in \WH(d)$ and $k \in \{0, \ldots, d-1\}$,
  \[
    \sigma_k\!\bigl(\phi(W)(\eps_d)\bigr) = (W v_d)(k),
  \]
  where $v_d \in \bC^d$ is defined by $v_d(k) = \sigma_k(\eps_d)$.
\end{lemma}

\begin{proof}
  By class field theory, $G_d \cong \mathrm{Cl}_{\mathfrak{m}_d}(K_d)$ via the Artin map. What this gives us is the abstract group structure of $G_d$. What it does not give us — and what we actually need — is a reason for $\WH(d)$ to map into $G_d$ at all. The Weyl-Heisenberg group is a quantum-mechanical object; the ray class group is an arithmetic one. That they are related is not a consequence of the Artin isomorphism; it is an independent conjecture, verified by Appleby et al.\ \cite{afmy2017} in every dimension where a SIC is numerically known, but not proved in general.

  Specifically: the claim is that $\WH(d)/Z(\WH(d)) \cong (\bZ/d\bZ)^2$ acts on $\Sigma_d$ via a subquotient of $\mathrm{Cl}_{\mathfrak{m}_d}(K_d)$, and that this Galois action on embeddings matches the WH action on $v_d$. We take this as H3. Condition (4) of Definition~\ref{def:frobenius-fixed} is the local shadow of this: it says that Frobenius elements at unramified primes act compatibly with WH displacements. Given that condition, the homomorphism
  \[
    \phi : \WH(d) \to \mathrm{Cl}_{\mathfrak{m}_d}(K_d) \xrightarrow{\mathrm{Artin}} G_d \hookrightarrow \mathrm{Aut}(L_d/\bQ)
  \]
  is well-defined and injective. But the work of showing condition (4) is satisfiable for a Frobenius-fixed $\eps_d$ is, in practice, where the argument would need to be filled in.
\end{proof}

This is the key structural move, and it carries real weight. The WH action on vectors in $\bC^d$ is broadcast simultaneously over the full Galois orbit — so equiangularity, once established, holds for all $d^2 - 1$ off-diagonal pairs at once. But establishing it requires the correspondence above, which is an open conjecture for general $d$.

\medskip\noindent\rule{\linewidth}{0.3pt}\medskip

\subsection{Simultaneous equiangularity}

The equiangularity argument requires one additional conjecture beyond H1--H3, which we state separately because its character is different from all three: it is a global harmonic analysis statement rather than a structural one.

\begin{quote}
\textbf{C4 (WH Irreducibility and Isotypic Isolation).} Under H2a--H2b and H3, the operators $\{\phi(W)\}_{W \in \WH(d)}$ form a commuting family acting on the Galois orbit $\{\sigma_k(\eps_d)\}$. C4 asserts three conditions simultaneously: (i) the WH representation on the Galois orbit $V_d = \mathrm{span}_\bC\{\sigma_k(\eps_d)\}$ is \emph{irreducible}; (ii) the isotypic component of $V_d$ corresponding to the trivial central character has \emph{multiplicity one}; and (iii) $v_d = (\sigma_0(\eps_d), \ldots, \sigma_{d-1}(\eps_d))$ is supported entirely in that component, so that
\[
  \langle v_d,\, \phi(W) v_d \rangle = \begin{cases} \|v_d\|^2 & W = \mathrm{Id}, \\ 0 & W \neq \mathrm{Id}. \end{cases}
\]
Conditions (i) and (ii) are necessary: without irreducibility, a joint eigenvector of the commuting algebra may have nonzero overlap with multiple isotypic components, which breaks the vanishing of off-diagonal WH matrix coefficients. ``Joint eigenvector of the commuting algebra'' is strictly weaker than ``vanishing of all off-diagonal WH Fourier coefficients''; it is conditions (i)--(ii) that bridge the two. After normalization, C4 yields $|\langle \psi_d, W\psi_d\rangle|^2 = 1/(d+1)$.

\end{quote}

C4 is a strong spectral purity condition: it demands that $v_d$ be spectrally pure for the WH action, not merely orthogonal in some averaged sense. In finite Fourier analysis terms, $v_d$ must correspond to a single WH character, which is equivalent to a Parseval identity on $G_d$-characters. The regular representation of $G_d$ is reducible, and lying in a single isotypic component is a significant constraint. Locally — at each unramified prime — a Hilbert~90 argument gives norm-1 elements as coboundaries, whose traces vanish. C4 asserts this holds globally; in the mixed-signature case that passage is not automatic. In the totally real and CM cases it follows from the Stark theorem. Here it may require either a strengthening of H1 to include an explicit global trace formula, or a separate id\`{e}lic argument.

\begin{lemma}[Equiangularity via Galois descent]
  \label{lem:equiangular}
  Assume H1, H2a, H3, and C4. After normalization $\psi_d = v_d / \|v_d\| \cdot \sqrt{d}$, we have $|\langle \psi_d, W\psi_d \rangle|^2 = 1/(d+1)$ for all non-identity $W \in \WH(d)$.
\end{lemma}

\begin{proof}
  By Lemma~\ref{lem:galois-weyl} (using H2a and H3), the inner product expands as
  \[
    \langle v_d, W v_d \rangle = \sum_{k=0}^{d-1} \sigma_k(\eps_d)\, \overline{\sigma_k(\phi(W)(\eps_d))}.
  \]
  Conditions H2a and H2b — specifically the palindromic minimal polynomial supplied by H2b — ensure that the product $\eps_d \cdot \phi(W)(\eps_d^{-1})$ lies in the kernel of the norm map $N_{L_d/K_d}$. Without palindromicity this fails. By Hilbert's Theorem~90, a norm-1 element in a Galois extension is a coboundary, and a coboundary's trace over the Galois orbit vanishes. For $W = \mathrm{Id}$ the argument gives $\|\eps_d\|^2$ instead. Conjecture C4 asserts this local coboundary conclusion holds globally over all embeddings simultaneously; given C4, the sum evaluates as stated and the modulus-squared condition follows after normalization.
\end{proof}

Uniqueness of the fiducial follows immediately from H2a: any other WH-covariant fiducial would produce a unit with identical Galois properties, differing from $\eps_d$ only by a root of unity, which corresponds to a global phase in Hilbert space.

\medskip\noindent\rule{\linewidth}{0.3pt}\medskip

\subsection{The main theorem}

\begin{theorem}[SIC-POVM existence, conditional]
  \label{thm:sic-existence}
  Assume H1, H2a, H3, and C4. Then for every $d \geq 1$, there exists a WH-covariant SIC-POVM in $\bC^d$.
\end{theorem}

\begin{proof}
  For $d = 1, 2$ the result is constructive. For $d \geq 3$ with $m_d$ non-square: H1 gives $\eps_d \in \OO_{L_d}^\times$; H2a and H2b give it the Frobenius symmetry of Definition~\ref{def:frobenius-fixed}. Define $v_d$ by $v_d(k) = \sigma_k(\eps_d)$, normalize to $\psi_d = v_d/\|v_d\| \cdot \sqrt{d}$, and for each $W \in \WH(d)$ let $\Pi_W = \frac{1}{d}(W\psi_d)(W\psi_d)^\dagger$.

  Lemma~\ref{lem:equiangular} (via H3, C4) gives $|\langle \psi_d, W\psi_d\rangle|^2 = 1/(d+1)$ for $W \neq \mathrm{Id}$. For $W \neq W'$, $|\langle W\psi_d, W'\psi_d\rangle|^2 = |\langle \psi_d, W^{-1}W'\psi_d\rangle|^2 = 1/(d+1)$. Equiangularity prevents collisions; the tight frame property follows by a standard argument \cite{appleby2019}. The case $m_d$ a perfect square is analogous with $K_d = \bQ$.
\end{proof}

H1, H2a, H3, and C4 are the four independent conjectures used; H2b is a derived consequence of H2a. All four conjectures remain open. The paper's claim is that they are \emph{sufficient}: given these four, the SIC-POVM construction is a theorem, with the frame property and uniqueness following cleanly. Whether identifying these four as the right targets constitutes progress is for Others to judge.

\section{Spectral Families}
\label{sec:families}

Since $K_d = \bQ(\sqrt{m_d})$ depends only on the square-free part of $m_d = d(d-2)$, all dimensions sharing the same square-free part share a base field. We call such a collection a \emph{spectral family}. The Stark Hypothesis is a statement about $K_d$, not about $d$ individually: to prove it for $K_d = \bQ(\sqrt{D})$, one proves the existence of Frobenius-fixed Stark units in the ray class fields over $\bQ(\sqrt{D})$. Once done, the result propagates to every dimension in the spectral family simultaneously.

\begin{corollary}
  \label{cor:spectral}
  If H1, H2a, H3, and C4 hold for all ray class fields over $K = \bQ(\sqrt{D})$, then WH-covariant SIC-POVMs exist for every dimension $d$ with $\mathrm{sqfree}(d(d-2)) = D$.
\end{corollary}

The corollary requires care. Sharing a base field $K_d$ does not automatically mean the SIC fields $L_d$ are compatible: the modulus $\mathfrak{m}_d$ varies with $d$, and two dimensions with the same $K_d$ may live in genuinely different ray class fields over it. H3 and C4 must hold for each $L_d$ separately. What the spectral family structure does give is that H1 — the Stark regulator statement — is at least asking about the same number field. Whether H2a propagates across the family is a separate question.

The observation that started this paper was that $d = 19$ and $d = 48$ have structurally identical SIC encodings in the grammar. We take that as a suggestion, not a theorem. Until the conductors $\mathfrak{m}_{19}$ and $\mathfrak{m}_{48}$ are verified against the Appleby et al.\ tables, the $19$/$48$ coupling is a numerical coincidence whose arithmetic content we have not established.

The more defensible general point: the variable $K_d$ captures something real about the difficulty of the SIC problem across dimensions. For each distinct base field there is one Stark-regulator question. That is a coarser invariant than the full set of conjectures, but it is a starting point for organizing which dimensions should be attacked together.

\begin{remark}[Field-level vs.\ ray-class-level invariants]
Two families of invariants govern the SIC problem over a spectral family, and they are controlled by different objects. The \emph{field-level} invariants are controlled by $K_d = \bQ(\sqrt{D})$: these include the unit group $\OO_{K_d}^\times$, the class number $h(K_d)$, and the Stark regulator (H1). They are shared across every dimension $d$ in the family. The \emph{ray-class-level} invariants are controlled by the modulus $\mathfrak{m}_d$: these include the ray class group $\mathrm{Cl}_{\mathfrak{m}_d}(K_d)$, the Galois group $G_d \cong \mathrm{Cl}_{\mathfrak{m}_d}(K_d)$, and the cohomological alignment of H3. Since $\mathfrak{m}_d$ depends on $d$, ray-class-level invariants are generally distinct across the family. In practical terms: H1 may be established once for $K_d$; H2a, H2b, H3, and C4 must each be verified for every $\mathfrak{m}_d$ separately.
\end{remark}

\section{Research Directions}
\label{sec:berry}

None of what follows is part of the proof architecture. We record it because it points toward how H1 and H2a might be attacked, and because the Berry phase identification seems like the kind of thing that should be written down even if we cannot yet make it rigorous.

The Stark conjecture, in its classical formulations, pins $\log|\eps_d|$ at each infinite place via the leading term of the partial zeta function at $s = 0$. What it does not fix is $\mathrm{Im}(\log \eps_d)$ — the imaginary part of the logarithm of the Stark unit. In the SIC context, this imaginary part has a direct physical interpretation: it is the Berry phase accumulated by the fiducial vector under the order-3 Zauner circuit.

To see this, note that the Zauner unitary $Z_d$ acts on $\psi_d$ by $Z_d \psi_d = e^{i\theta_d} \psi_d$ for some phase $\theta_d$ (since $\psi_d$ is a fixed point of $Z_d$ by condition (1) of Definition~\ref{def:frobenius-fixed}). Under the Galois-to-Weyl correspondence, this phase is
\[
  \theta_d = \mathrm{Im}(\log \eps_d).
\]
The Weyl-Heisenberg equiangularity condition constrains $\theta_d$ through the requirement that the inner products $\langle \psi_d, W\psi_d \rangle$ have modulus $1/\sqrt{d+1}$. If the equiangularity condition pins $\theta_d$ uniquely, then together with the Stark regulator data (which fixes $|\eps_d|$), the unit $\eps_d$ is fully determined. Proving existence would then reduce to a compatibility problem between the real Stark data (which fixes $|\eps_d|$ via the regulator) and the imaginary Stark data (which would need to supply $\mathrm{Im}(\log \eps_d)$ consistent with the equiangularity constraint).

We do not know whether the equiangularity condition uniquely pins $\theta_d$. It does in the dimensions we have checked, but we have not found a general argument. If true, it would mean that the SIC-POVM existence conjecture and the imaginary part of the mixed-signature Stark conjecture are equivalent — a reduction that would make both problems accessible from either direction.

We also do not know whether to connect this to the Langlands program. The Galois-Weyl correspondence of Lemma~\ref{lem:galois-weyl} has the rough shape of a Langlands functoriality, with the ray class field $L_d$ on the Galois side and the Weyl-Heisenberg orbit in $\bC^d$ on the automorphic side. We find this analogy compelling but have not been able to make it precise enough to be useful.

\section{Open Problems}

We list these not as a checklist but as an attempt to be honest about exactly where the argument is incomplete.

\begin{enumerate}
  \item \textbf{Prove H2a.} The mixed-signature Stark conjecture (H1) is the most classical of the four open problems, and in some sense the most tractable — it fits into a framework that already has a proof strategy in the totally real and CM cases. H2a is harder to locate in the literature, and we think it is the true crux. The question is whether the $(-1)$-eigenspace of complex conjugation on the Dirichlet regulator matrix is large enough to contain a generator. We do not know a reason it should fail, but we also do not know a reason it should hold, and that asymmetry is uncomfortable.

  \item \textbf{Prove H3.} The claim that $\WH(d)$ acts on the Galois embeddings of $L_d$ via a subquotient of $\mathrm{Cl}_{\mathfrak{m}_d}(K_d)$ is numerically confirmed in every dimension where a SIC is known, but the question of why it should be true has not had a satisfying answer. We now think it does.

  The correspondence, we believe, is the restriction of geometric Langlands functoriality to the abelian sector for $K_d = \bQ(\sqrt{d(d-2)})$ — specifically, the abelian case of the functoriality for real quadratic fields. What would need to be proved is a Uniformization Theorem identifying $\WH(d)/Z(\WH(d)) \cong (\bZ/d\bZ)^2$ as a geometric boundary encoding of the arithmetic bulk of $\mathrm{Cl}_{\mathfrak{m}_d}(K_d)$: the quantum phase space group as a imscriptive boundary whose bulk is the class group. In this framing, the correspondence is not a coincidence but a consequence of the abelian case of the Langlands dictionary — closer in character to Hilbert's 12th problem than to the Stark conjecture.

  H1 and H3 are independent. H1 is about existence and regulator properties of $\eps_d$; it says nothing about how $\WH(d)$ relates to $\mathrm{Cl}_{\mathfrak{m}_d}(K_d)$. Anyone working toward the main theorem needs to treat them as separate targets with separate techniques.

  \item \textbf{Control the global trace sum.} We noted in §\ref{sec:proof} that an earlier attempt to use Chebotarev density to close the equiangularity argument fails: local coboundaries do not assemble into a global identity without additional input. This is embarrassing — it is the kind of gap that looks as if it should close by a standard argument. And yet, we have not found one. Several colleagues have also looked. They have not found one either, though a great many reported that they had, in fact, identified better uses of their time. Resolving it may require strengthening H1 to include an explicit global trace formula, or a separate id\`{e}lic argument.

  \item \textbf{Verify the $d = 19$, $d = 48$ coupling.} The observation that started this paper deserves to be confirmed. Check the conductors $\mathfrak{m}_{19}$ and $\mathfrak{m}_{48}$ against the Appleby et al.\ tables; until that is done, the coupling is a numerical suggestion, not a theorem.

  \item \textbf{The Berry phase equivalence.} Determine whether the equiangularity condition uniquely pins $\mathrm{Im}(\log \eps_d)$. The stakes are higher than they may appear: if it does, the SIC conjecture and the imaginary Stark conjecture are equivalent, making both problems simultaneously accessible from whichever direction turns out to be more tractable.

  \item \textbf{Spectral family census.} For each square-free $D$, identify all $d$ with $\mathrm{sqfree}(d(d-2)) = D$. The families containing currently unproven dimensions are the natural targets.
\end{enumerate}

\section*{Conclusion}

We set out to explain a numerical coincidence — identical structural encodings for $d = 19$ and $d = 48$ — and ended up with what we think is a more honest account of why the SIC conjecture has been hard. The difficulty is not computational. It is that the equiangularity condition requires an exact $Z_2$ self-duality in the underlying unit group, and that duality is not something you can approximate your way into: the non-synthesizability theorem says it must be present from the beginning, in the object you start with.

The Stark unit is an object that starts with it. Its palindromic self-duality is not imposed from the outside but intrinsic to its Galois symmetry profile. In that sense, the palindromicity of the minimal polynomial and the equiangularity of the SIC are not analogous structures that happen to match — they are the same structure, read off from the same algebraic datum in two different coordinate systems.

We have not proved that this Stark unit exists with all the properties we need. That is the paper's honest limitation, and we have tried not to bury it. What we have done is reduce the existence question to four identifiable conjectures — one classical, two new, one harmonic-analytic — and show that once those four hold, nothing else is required. Whether that reduction constitutes progress, and whether the four conjectures are more accessible than the original, others will have to judge. We believe the answer is yes on both counts.

\vspace{1em}
\noindent\textbf{Acknowledgments.} We thank the Appleby school for making their conductor data and fiducial coordinates publicly available, without which the spectral family observation would not have been possible. We would also like to thank Harry T. Larson and his insistence that we see out a lemma to its ends. This work was supported by the author's independent research.

\section*{Declarations}

\noindent\textbf{Competing Interests.} The author has no relevant financial or non-financial interests to disclose.

\noindent\textbf{Data Availability.} This manuscript presents a mathematical proof and does not rely on or generate empirical data. No datasets were created or analyzed during this study.

\noindent\textbf{Funding.} This work was supported by the author's independent research program.

\begin{thebibliography}{9}

\bibitem{zauner1999}
Zauner, G.: Quantendesigns: Grundz\"uge einer nichtkommutativen Designtheorie. Ph.D.\ thesis, Universit\"at Wien (1999)

\bibitem{afmy2017}
Appleby, M., Flammia, S., McConnell, G., Yard, J.: SICs and algebraic number theory. Found.\ Phys.\ \textbf{47}, 1042--1059 (2017). \url{https://doi.org/10.1007/s10701-017-0082-9}

\bibitem{synthref}
Mills, L.: SynthOmnicon (2026). \url{https://github.com/umpolungfish/synthomnicon}

\bibitem{stark1971}
Stark, H.M.: Values of $L$-functions at $s=1$. Adv.\ Math.\ \textbf{7}, 301--343 (1971)

\bibitem{tate1984}
Tate, J.: Les Conjectures de Stark sur les Fonctions $L$ d'Artin en $s=0$. Birkh\"auser, Basel (1984)

\bibitem{appleby2019}
Appleby, M., Bengtsson, I., Flammia, S., Goyeneche, D.: Tight frames, Hadamard matrices and Zauner's conjecture. J.\ Phys.\ A \textbf{52}, 295301 (2019). \url{https://doi.org/10.1088/1751-8121/ab25ad}

\end{thebibliography}

\end{document}
